%%%%%%%%%%%%%%%%%%%%%%%%%%%%%%%%%
%
%       EXERCÍCIO
%
%%%%%%%%%%%%%%%%%%%%%%%%%%%%%%%%%

\definevalorquestao

\begin{exercicioBanco}[\valorquestao]
O custo para produzir um produto é \(C(x) = 50 + 2x + 0,1x^{2}\), onde \(x\) é a quantidade produzida do produto. Cada unidade é vendida por R\$\(6,50\). Entre que valores deve variar \(x\) para não haver prejuízo?

% Define as alternativas
\newcommand{\alternativas}{%
    \begin{center}
        \begin{tabularx}{\textwidth}{XXX}
            (a) \(x \le 20\) ou \(x \ge 25\). &
            (b) \(10 \le x \le 50\). &
            (c) \(x \ge 20\). \\[5pt]
            (d) \(x \le 25\). &
            (e) \(20 \le x \le 25\).
        \end{tabularx}
    \end{center}
}

% Define a resposta correta
\newcommand{\resposta}{E}

% Lógica condicional para exibição
\ifdefstring{\atividade}{lista}{%
    \alternativas
    \vspace{0.5em}
    
    \noindent\textbf{Resposta:} letra \textbf{\resposta}.
}{%
    \ifdefstring{\modo}{objetivo}{%
        \alternativas
    }{%
        % modo = discursiva → não mostra alternativas
    }
}
\end{exercicioBanco}

